\chapter{Notación y preliminares}
\label{ch:notation}

Este capítulo fija la notación usada en todo el documento. Las
convenciones son compatibles con CLRS donde aplica.

\section{Grafos}

\begin{definition}[Multigrafo dirigido temporal]
Un \emph{multigrafo dirigido temporal} es una tupla $G = (V, E,
\tau, \ell_V, \ell_E)$ donde $V$ es un conjunto finito de
vértices, $E \subseteq V \times V \times \mathbb{N}$ es un
multiconjunto de aristas dirigidas (el tercer componente es un
identificador único de arista, por lo que varias aristas entre el
mismo par son elementos distintos de $E$), $\tau : E \to
\mathbb{Z}$ asigna a cada arista una marca temporal de evento,
$\ell_V : V \to \mathcal{T}_V$ etiqueta a los vértices con un tipo
de entidad, y $\ell_E : E \to \mathcal{T}_E$ etiqueta a las
aristas con un tipo de relación.
\end{definition}

Escribiremos $|V| = n$ y $|E| = m$. La vecindad saliente de un
vértice $v \in V$ en la ventana $[t_0, t_1]$ es
\[
    N^{+}_{[t_0, t_1]}(v) = \{\, u \in V \mid \exists\, e = (v, u,
    k) \in E : \tau(e) \in [t_0, t_1]\,\}.
\]
La vecindad entrante $N^{-}_{[t_0, t_1]}(v)$ se define
simétricamente. Cuando la ventana es todo el dominio omitimos el
subíndice y escribimos $N^{+}(v)$, $N^{-}(v)$.

\begin{definition}[Extensión MVCC]
La \emph{extensión MVCC} de $G$ agrega un segundo timestamp
$\iota : E \to \mathbb{Z}$ que registra cuándo se insertó cada
arista en el grafo. El motor almacena $\iota$ como un delta
compacto contra un epoch fijo; ver
Capítulo~\ref{ch:compact-relation}. Una consulta por snapshot en
el instante $s$ considera sólo el subgrafo
\[
    G|_{s} = (V, \{\,e \in E \mid \iota(e) \leq s\,\}, \tau, \ell_V,
    \ell_E),
\]
es decir, ``lo que el grafo sabía al momento $s$''.
\end{definition}

\section{Caminos e hipótesis}

\begin{definition}[Camino temporal]
Un \emph{camino temporal} de longitud $\ell$ en $G$ es una
secuencia
\[
    \pi = (v_0, e_1, v_1, e_2, \ldots, e_\ell, v_\ell),
\]
con $v_i \in V$, $e_i = (v_{i-1}, v_i, k_i) \in E$, y
$\tau(e_1) \leq \tau(e_2) \leq \cdots \leq \tau(e_\ell)$
(monotonía temporal). Con frecuencia abreviaremos un camino como
su lista de vértices $(v_0, v_1, \ldots, v_\ell)$ cuando las
aristas intermedias se deducen del contexto.
\end{definition}

\begin{definition}[Camino $k$-simple]
Un camino es \emph{$k$-simple} si y sólo si ningún vértice aparece
en él más de $k$ veces. La $k$-simplicidad relaja la exigencia
clásica de no-repeticiones ($k=1$) para admitir ciclos que un
atacante real sí exhibe --- por ejemplo, un bucle de callback C2
donde el mismo proceso vuelve a entrar en el camino. $k$ se fija
por hipótesis; el valor por defecto es 1.
\end{definition}

\begin{definition}[Hipótesis]
Una \emph{hipótesis} $H$ es una cadena de flechas
\[
    H = (\sigma_0, \rho_1, P_1^{\text{edge}}, P_1^{\text{dst}},
        \sigma_1, \rho_2, \ldots, \rho_\ell, P_\ell^{\text{edge}},
        P_\ell^{\text{dst}}, \sigma_\ell)
\]
donde $\sigma_i \in \mathcal{T}_V \cup \{\ast\}$ es el tipo de
entidad requerido en el $i$-ésimo vértice (o $\ast$ para
``cualquiera''), $\rho_i \in \mathcal{T}_E \cup \{\ast\}$ es el
tipo de relación requerido en la $i$-ésima arista, y
$P_i^{\text{edge}}$, $P_i^{\text{dst}}$ son predicados opcionales
sobre los metadatos de la arista y del vértice de destino,
respectivamente. La cadena también lleva restricciones de
agregación (\code{HAVING}, \code{WITHIN}) que se procesan después
del match. Un camino $\pi = (v_0, \ldots, v_\ell)$
\emph{satisface} $H$ si y sólo si $\ell_V(v_i) \in \sigma_i$,
$\ell_E(e_i) \in \rho_i$, y cada predicado se cumple sobre la
arista o el vértice correspondiente.
\end{definition}

\section{Notación asintótica}

Valen las convenciones estándar de CLRS. $\Ord{f(n)}$ es cota
superior, $\Thet{f(n)}$ es cota ajustada, $\Omega(f(n))$ es cota
inferior. ``Amortizado $O(f)$'' significa que el costo por
operación, promediado sobre una secuencia peor caso, es $O(f)$,
aun cuando operaciones individuales sean más caras.

\section{Modelo de memoria}

El modelo de propiedad (\emph{ownership}) de Rust hace que la
mayoría de los invariantes sean chequeables localmente. Donde el
motor depende de mutabilidad interior, lo marcamos explícitamente
--- \code{Arc<RwLock<\_>>} para estado de sesión compartido,
\code{RwLock<SpillableEdgeStore>} para el buffer de aristas en
memoria que puede volcarse a disco durante ingesta,
\code{OnceLock} para el snapshot del catálogo a nivel proceso. El
motor nunca comparte grafos mutables entre threads sin un lock; el
matcher DFS paraleliza dando a cada worker rayon una referencia
inmutable y una pila scratch local al thread
(Capítulo~\ref{ch:matcher}).

\section{Convención de punteros a fuente}

Cada capítulo termina con un bloque \emph{En el código} que lista
punteros \code{archivo:línea} a la implementación Rust. Las rutas
son relativas a la raíz del repositorio, p.~ej.\
\filepath{graph\_hunter\_core/src/graph.rs:461}. Los números de
línea pueden derivar entre revisiones, así que si el puntero está
desactualizado conviene grepear el nombre de la función.
